% =====================================================================
%  Chapter 16 — The map
% =====================================================================

\chapter{The Method in the Source: A Step-to-Declaration Map}
\label{ch:map}

This chapter is the apparatus that makes the preceding ten chapters
checkable. For each step of the method it gives the subtree of the submitted
development where the mathematics lives, and the specific declarations a reader
should look at.

Everything here was read from the submitted source tree at the pinned
revision. Paths are relative to the archive root.

\section{Where to start if you want to check one step}\begin{center}
\small
\setlength{\extrarowheight}{-1pt}
\begin{tabular}{@{}>{\RaggedRight\arraybackslash\hyphenpenalty=10000\exhyphenpenalty=10000}p{0.22\textwidth}>{\RaggedRight\arraybackslash\hyphenpenalty=10000\exhyphenpenalty=10000}p{0.37\textwidth}>{\RaggedRight\arraybackslash\hyphenpenalty=10000\exhyphenpenalty=10000}p{0.37\textwidth}@{}}
\toprule
\textbf{If you want to check} & \textbf{Go to} & \textbf{And read} \\
\midrule
the statement of the claim & \code{src/Statements.lean} &
\code{RiemannZetaZerosOnLine}, \code{DirichletLNonvanishingRightOfHalf} \\
the reduction to RH & \code{CharacterBoxes/RiemannHypothesisOfEveryCharacter.lean} &
\code{riemannHypothesis\_of\_everyCharacter} \\
the theta lift & \code{.../ThetaParseval/HalfIndexLattices/TLiftK.lean} &
\code{UpperTheory.TP.TLiftK} \\
the chain and its cost & \code{Formalism/ThetaLiftIndependence/Statement.lean} &
\code{KerComp}, \code{ChainInt}, \code{ChainIntM} \\
the two thirds & \code{Formalism/SharpRestriction/GrowthBounds/SegmentDegree/} &
\code{ThresholdPowerBoundOfCounts.lean}, \code{AtypicalPathSmallWeight.lean} \\
the seesaw & \code{.../Divisibility/FunctionModelRallisFromInputs.lean} &
\code{bB}, \code{vth}, \code{PerNZ}, \code{Conclusion1} \\
the averaging & \code{Formalism/BehrendSets/Statement.lean} &
\code{BehrendFree}, \code{BehrendBound}, \code{TranslateMean} \\
the vanishing & \code{AbelianVarietyTheory/Definitions/Foundations/FinalTheorems.lean} &
\code{admitted}, \code{VanishOrd} \\
the bottom of the chain & \code{.../SatakeParameters/LefschetzFixedPointFiniteEtaleCorrespondence.lean} &
\code{lefschetz\_fixedPoint\_finite\_etale\_correspondence} \\
\bottomrule
\end{tabular}
\end{center}

\section{The full map}\begin{longtable}{@{}>{\RaggedRight\arraybackslash\hyphenpenalty=10000\exhyphenpenalty=10000}p{0.13\textwidth}>{\RaggedRight\arraybackslash\hyphenpenalty=10000\exhyphenpenalty=10000}p{0.40\textwidth}>{\RaggedRight\arraybackslash\hyphenpenalty=10000\exhyphenpenalty=10000}p{0.41\textwidth}@{}}
\toprule
\textbf{Step} & \textbf{Subtree} & \textbf{Key declarations} \\
\midrule
\endhead
\bottomrule
\endfoot
1 & \code{Development/WeilBound/RepresentationTheoryGL/SiegelWeilProof/} (635
files), \code{.../ThetaLifts/LeviDeterminants/} (116) &
\code{FrameBasis}, \code{ExistsColsDetNeZero}, \code{DomWeil},
\code{sigRapid\_of}, \code{cuspRapid\_of} \\
2 & \code{Formalism/SharpRestriction/Divisibility/},
\code{.../Divisibility/KudlaSiegelWeil.lean} &
\code{AIntO}, \code{projector\_statement\_of\_integrality},
\code{thP}, \code{pow\_two\_pow\_add\_inv\_succ},
\code{pow\_two\_pow\_add\_inv\_mem\_adjoin}, \code{exists\_sqrt\_two} \\
3 & \code{Development/WeilBound/RepresentationTheoryGL/RegularisedSiegelWeil/}
(153), \code{.../CoreStatements/HomogeneousNonvanishing/} (105),
\code{AbelianVarietyTheory/Proofs/DecoratedLatticeProofs/} &
\code{weilCoefficient\_splits\_localFactor\_twistPlace},
\code{weilCoefficient\_pureTensor\_eq\_smul\_inner},
\code{weilCoefficient\_externalProduct\_eq\_mul\_orthogonalSum},
\code{GramMatrixTwistedLieModels} \\
4 & \code{Formalism/ThetaLiftIndependence/},
\code{Development/WeilBound/ThetaLifts/} &
\code{chainOp}, \code{chainOpM}, \code{KerComp}, \code{KerAdjoint},
\code{ChainInt}, \code{ChainIntM}, \code{OrbitSum} \\
5 & \code{Formalism/SharpRestriction/GrowthBounds/SegmentDegree/} (67 files),
\code{.../Divisibility/ChainOperators/}, \code{.../Divisibility/CutKapD.lean},
\code{.../Divisibility/CutDivD.lean} &
\code{threshold\_arith\_letter\_case}, \code{threshold\_arith\_w\_case},
\code{threshold\_power\_bound\_of\_counts}, \code{Typ}, \code{GoodA},
\code{order\_numeric\_bound\_of\_oldest}, \code{rank\_input\_holds},
\code{kappa\_degree\_bound\_holds}, \code{sum\_char},
\code{LocalAdditiveCharacterExpansion}, \code{projector\_statement\_of\_stability} \\
6 & \code{.../Divisibility/FunctionModelRallisFromInputs.lean},
\code{.../DecoratedLatticeProofs/ThetaInnerProductFormulaDefiniteOfRallisConverge.lean} &
\code{bB}, \code{vth}, \code{PerNZ}, \code{IsAutRep}, \code{IsDv}, \code{IsMu},
\code{Hyp3}, \code{TParity}, \code{NCGaps}, \code{car}, \code{cenY},
\code{ThetaDescentTowerConditions}, \code{Conclusion1} \\
7 & \code{Formalism/BehrendSets/}, \code{.../BehrendVariance/} &
\code{BehrendFree}, \code{BehrendBound}, \code{ThickSet}, \code{TranslateMean},
\code{OrbitOne}, \code{WindowTwo}, \code{SumSq}, \code{SumLin},
\code{ShortOrdLarge}, \code{ShortOrdSq}, \code{ShortOrdLin}, \code{OrdGrp},
\code{frob}, \code{OrdCard}, \code{OrdIff} \\
8 & \code{Development/AutomorphicArithmetic/AutomorphicVarieties/TheoremsMuLinked/},
\code{Development/WeilBound/CoreStatements/AutomorphicKernels/} (932),
\code{.../CoreStatements/PublishedTheorems/} (540) &
\code{GlobalTheory.Mem.LiteralAutomorphy.LitN}, the Galois-representation
types in \code{TheoremsMuLinked/GaloisRepresentationTypes/}, and the
Poincar\'e-series machinery in \code{AutomorphicKernels/} \\
9 & \code{AbelianVarietyTheory/Definitions/Foundations/FinalTheorems.lean},
\code{.../Proofs/GlobalInterpolation/NonemptyHypsOfAllRows.lean} &
\code{admitted}, \code{VanishOrd}, \code{Nf}, \code{Unr}, \code{Gr},
\code{Triv}, \code{GrZero}, \code{nonempty\_hyps\_of\_allRows} \\
\bottomrule
\end{longtable}

\section{The chain from the headline theorem to the bottom of the mathematics}This is the path traced in Chapter~\ref{ch:chain}, collected in one place.

\begin{longtable}{@{}>{\RaggedRight\arraybackslash\hyphenpenalty=10000\exhyphenpenalty=10000}p{0.36\textwidth}>{\RaggedRight\arraybackslash\hyphenpenalty=10000\exhyphenpenalty=10000}p{0.56\textwidth}@{}}
\toprule
\textbf{Declaration} & \textbf{File} \\
\midrule
\endhead
\bottomrule
\endfoot
\code{riemannHypothesis} & \code{src/Main.lean} \\
\code{UpperTheory.FinalTheorems.riemannHypothesis} & \code{CharacterBoxes/MainConclusion.lean} \\
\code{UpperTheory.riemannHypothesis\_of\_everyCharacter} & \code{CharacterBoxes/RiemannHypothesisOfEveryCharacter.lean} \\
\code{UpperTheory.EveryCharacter} & \code{CharacterBoxes/EveryCharacter.lean} \\
\code{UpperTheory.FinalTheorems.everyCharacter} & \code{CharacterBoxes/MainConclusion.lean} \\
\code{UpperTheory.Gal.Int.galoisLocalGlobalStrong\_recParameter} & \code{.../IntegralStructures/RowSrc.lean} \\
\code{UpperTheory.Gal.Int.galoisRepresentations\_of\_thetaLift} & \code{.../IntegralStructures/GaloisRepresentationsOfThetaLift.lean} \\
\code{UpperTheory.Gal.Int.galoisLocalGlobalStrong\_of\_geometry\_lift} & \code{.../IntegralStructures/Development.lean} \\
\code{UpperTheory.Gal.Int.Geo.geoOne\_holds} & \code{.../GeometricFlagBases/WeightOneGeometricRealization.lean} \\
\code{UpperTheory.Gal.Int.Geo.weight\_one\_geometric\_realization} & same file \\
\code{GlobalTheory.T.lefschetz\_fixedPoint\_finite\_etale\_correspondence} & \code{.../SatakeParameters/LefschetzFixedPointFiniteEtaleCorrespondence.lean} \\
\code{UpperTheory.Cb.lefschetz\_fixedPoint\_finite\_etale\_correspondence} & \code{.../EtaleClopenDescent/LefschetzFixedPointFiniteEtaleCorrespondence.lean} \\
\bottomrule
\end{longtable}

\section{The trees, by size}A reader who wants to understand the shape of the argument before reading any of
it can start with the sizes, which were computed from the submission.

\begin{center}
\small
\setlength{\extrarowheight}{-1pt}
\begin{tabular}{@{}>{\RaggedRight\arraybackslash\hyphenpenalty=10000\exhyphenpenalty=10000}p{0.34\textwidth}>{\RaggedRight\arraybackslash\hyphenpenalty=10000\exhyphenpenalty=10000}p{0.12\textwidth}>{\RaggedRight\arraybackslash\hyphenpenalty=10000\exhyphenpenalty=10000}p{0.50\textwidth}@{}}
\toprule
\textbf{Directory} & \textbf{Files} & \textbf{What it holds} \\
\midrule
\code{Development/} & 141{,}928 & the theorems themselves \\
\quad \code{WeilBound/} & --- & the Weil bound: Siegel--Weil, theta lifts,
frame estimates \\
\quad \quad \code{CoreStatements/AutomorphicKernels/} & 932 & automorphic
kernels, Poincar\'e series \\
\quad \quad \code{CoreStatements/PublishedTheorems/} & 540 & the published
results the development imports \\
\quad \quad \code{RepresentationTheoryGL/SiegelWeilProof/} & 635 & the
Siegel--Weil proof \\
\quad \quad \code{RepresentationTheoryGL/RegularisedSiegelWeil/} & 153 & the
regularised form \\
\quad \quad \code{ThetaLifts/} & --- & the lifts and their estimates \\
\quad \code{AutomorphicVarieties/TheoremsMuLinked/} & --- & automorphy of
$\mu$ \\
\code{Formalism/} & 18{,}254 & the general machinery \\
\quad \code{SharpRestriction/Divisibility/} & --- & divisibility, the chain,
the budget \\
\quad \code{SharpRestriction/GrowthBounds/SegmentDegree/} & 67 & the segment
degree bounds --- \emph{where the two thirds is} \\
\quad \code{BehrendSets/} & --- & Behrend's construction \\
\quad \code{ThetaLiftIndependence/} & --- & the chain operators \\
\bottomrule
\end{tabular}
\end{center}

\section{What the map does and does not license}\begin{quote}
\textbf{It does not license a verdict.} The map tells a reader where the
mathematics is. It does not tell them that the mathematics is right, and it
cannot: the statements in these files are the author's re-derivations of
published results, and whether each re-derivation is faithful is exactly the
question that has not been answered.
\end{quote}

What the map does license is a much sharper form of the question. Before the
map, the question was ``is this proof of RH?''. After it, the question is:
\emph{are these forty-odd declarations faithful restatements of the published
results they are named for, and are they applied in settings where those
results hold?} That is a finite, well-defined, expert-sized question --- and
Chapter~\ref{ch:protocol} gives the order to ask it in.

\begin{keytakeaways}
\item Every step of the method maps to a specific subtree, and every step maps
to specific named declarations. A reader can check any one step in isolation.
\item The two thirds is in \code{GrowthBounds/SegmentDegree/}, 67 files, in the
threshold arithmetic of \code{ThresholdPowerBoundOfCounts.lean}.
\item The full chain from the headline theorem to the proved Lefschetz formula
is eleven declarations in eight files.
\item The map sharpens the open question from ``is this a proof?'' to ``are
these declarations faithful, and are they applied where the classical results
hold?'' --- a finite question with a defined answer.
\end{keytakeaways}
